\documentclass[12pt]{article}

% --------------------
% PACOTES
% --------------------
\usepackage[utf8]{inputenc}
\usepackage[T1]{fontenc}
\usepackage[brazil]{babel}

\usepackage{amsmath, amssymb}
\usepackage{graphicx}
\usepackage{geometry}
\usepackage{hyperref}
\usepackage{booktabs}

\geometry{margin=2.5cm}

% --------------------
% COMANDOS
% --------------------
\newcommand{\ema}{Estações Meteorológicas Automáticas}
\newcommand{\mqtt}{MQTT}
\newcommand{\api}{API}

% --------------------
% DOCUMENTO
% --------------------
\title{Plano de Pesquisa \\ Sistema de \ema}
\author{Henry Campos}
\date{\today}

\begin{document}

\maketitle
\tableofcontents
\newpage

% --------------------
\section{Introdução}

O projeto propõe o desenvolvimento de um sistema de \ema\ distribuídas na cidade de Três Lagoas, com coleta e processamento de dados climáticos em tempo real.

% --------------------
\section{Objetivos}

\subsection{Objetivo Geral}
Desenvolver uma infraestrutura completa para monitoramento meteorológico urbano.

\subsection{Objetivos Específicos}
\begin{itemize}
    \item Implementar estações automáticas
    \item Criar backend com \api
    \item Desenvolver visualização web e mobile
\end{itemize}

% --------------------
\section{Metodologia}

O sistema será baseado em comunicação via \mqtt, com armazenamento em banco de dados temporal.

% --------------------
\section{Cronograma}

\begin{itemize}
    \item Etapa 1: Modelagem
    \item Etapa 2: Implementação
    \item Etapa 3: Testes
\end{itemize}

% --------------------
\bibliographystyle{plain}
\bibliography{referencias}

\end{document}
